% !TeX root = main.tex

Datenbanktransaktionen können aufgrund von Fehlern wie Deadlocks, Serialization Failures und Lock Timeouts fehlschlagen.
Retry-Mechanismen ermöglichen eine automatische Wiederholung fehlgeschlagener Transaktionen und können dadurch die
Erfolgsrate erhöhen.
Gleichzeitig verursachen zusätzliche Wiederholungsversuche weitere Ausführungszeit und Datenbanklast.
Diese Arbeit untersucht daher den Einfluss verschiedener Retry-Strategien und deren Parameter auf die Performance und
Zuverlässigkeit von PostgreSQL bei den genannten Fehlerklassen.
Hierzu wurden verschiedene Retry-Strategien sowie die Parameter Retry-Count und Delay unter unterschiedlichen
Concurrency-Werten systematisch untersucht.
Die Bewertung erfolgte anhand der Erfolgsrate und der Ausführungszeit.
Ergänzend wurden Retry-Verteilung, Quantile, Overhead und Durchsatz betrachtet.
Die Ergebnisse zeigen, dass Retry die Erfolgsrate der untersuchten Fehlerklassen grundsätzlich erhöhen kann,
wobei der Nutzen und die entstehenden Performancekosten stark von der jeweiligen Fehlerklasse und Konfiguration
abhängen.
Eine höhere Anzahl an Wiederholungsversuchen führt nicht grundsätzlich zu einem besseren Ergebnis, da zusätzliche
Ausführungszeit entsteht.
Auch der Delay beeinflusst das Verhältnis zwischen Erfolgsrate und Performance.
Für die untersuchten Fehlerklassen konnten unterschiedliche geeignete Konfigurationen identifiziert werden,
sodass keine allgemein optimale Retry-Konfiguration abgeleitet werden kann.
Die Untersuchung zeigt damit, dass Retry-Mechanismen grundsätzlich zur Erhöhung der Zuverlässigkeit beitragen können,
ihre Strategie und Parameter jedoch an die jeweilige Fehlerklasse und die Systembedingungen angepasst werden müssen.
